\section{Parte 1: Programação Estruturada e Operadores}

% Slide de Abertura da Secao 1
\begin{frame}
    \begin{center}
        \LARGE \textbf{Parte 1: Programação Estruturada}\\
        \vspace{0.3cm}
        \normalsize \textit{Operadores Aritméticos, Relacionais e Lógicos}\\
        \vspace{0.4cm}
        \small Questões 01 a 05
    \end{center}
    \vspace{0.3cm}
    \begin{alertblock}{Conceito Central}
        A programação estruturada baseia-se no fluxo sequencial e na manipulação de dados na memória através de expressões matemáticas e booleanas.
    \end{alertblock}
\end{frame}

% Questao 1
\begin{frame}{Questão 01: Operadores Aritméticos Básicos}
    \begin{block}{Quesito: Custo de Deslocamento e Consumo}
        Desenvolva um programa que receba três valores:
        \begin{enumerate}
            \item Distância total da viagem em quilômetros (\texttt{km});
            \item Consumo médio do veículo em quilômetros por litro (\texttt{km/l});
            \item Preço do litro de combustível em reais (\texttt{R\$}).
        \end{enumerate}
        O programa deve calcular e exibir:
        \begin{itemize}
            \item O total de litros consumidos ($\text{Distância} / \text{Consumo}$);
            \item O custo total em reais formatado com 2 casas decimais.
        \end{itemize}
    \end{block}

    \begin{exampleblock}{Ajuda de Resolução (Como Pensar)}
        \begin{itemize}
            \item Converta os dados lidos do teclado para \texttt{float}.
            \item Calcule os litros com o operador de divisão (\texttt{/}).
            \item Obtenha o custo multiplicando litros pelo preço unitário (\texttt{*}).
            \item Exiba o resultado com \texttt{f"\{custo:.2f\}"}.
        \end{itemize}
    \end{exampleblock}
\end{frame}

% Questao 2
\begin{frame}{Questão 02: Divisão Inteira e Operador Módulo}
    \begin{block}{Quesito: Conversão de Segundos para H:M:S}
        Um monitor de infraestrutura registra o tempo de atividade de um servidor em segundos inteiros. \\
        Escreva um programa que leia esse valor inteiro e o decomponha no formato: \textbf{Horas}, \textbf{Minutos} e \textbf{Segundos restantes}.
    \end{block}

    \begin{exampleblock}{Ajuda de Resolução (Como Pensar)}
        \begin{itemize}
            \item Lembre-se: $1\text{ h} = 3600\text{ s}$ e $1\text{ min} = 60\text{ s}$.
            \item Calcule as horas completas com divisão inteira: $\text{horas} = \text{total} // 3600$.
            \item Extraia o que sobrou com o operador resto (\texttt{\%}): $\text{resto} = \text{total} \% 3600$.
            \item Obtenha os minutos com $\text{minutos} = \text{resto} // 60$ e os segundos finais com $\text{segundos} = \text{resto} \% 60$.
        \end{itemize}
    \end{exampleblock}
\end{frame}

% Questao 3
\begin{frame}{Questão 03: Operadores Relacionais (Comparação)}
    \begin{block}{Quesito: Monitor de Capacidade de Link de Rede}
        Um link opera a $1000\text{ Mbps}$. Leia o tráfego atual (\texttt{Mbps}) e gere 3 variáveis booleanas (\texttt{True/False}), sem usar \texttt{if}:
        \begin{itemize}
            \item \textbf{Uso Seguro:} Tráfego estritamente menor que $80\%$ da capacidade;
            \item \textbf{Atenção Crítica:} Tráfego atingiu ou superou $90\%$ da capacidade;
            \item \textbf{Sobrecarga:} Tráfego estritamente superior a $100\%$ da capacidade.
        \end{itemize}
        Exiba o resultado booleano de cada uma das três avaliações.
    \end{block}

    \begin{exampleblock}{Ajuda de Resolução (Como Pensar)}
        \begin{itemize}
            \item Operadores de comparação (\texttt{<}, \texttt{>=}, \texttt{>}) retornam diretamente \texttt{True} ou \texttt{False}.
            \item Calcule os limites multiplicando a capacidade por $0.80$, $0.90$ e $1.0$.
            \item Guarde cada comparação direto em uma variável booleana e a imprima.
        \end{itemize}
    \end{exampleblock}
\end{frame}

% Questao 4
\begin{frame}{Questão 04: Operadores Lógicos Comportamentais}
    \begin{block}{Quesito: Controle de Acesso ao Datacenter}
        A catraca eletrônica deve autorizar a abertura da porta quando:
        \begin{itemize}
            \item O usuário possui credencial (\texttt{True/False}) \textbf{E} o horário atual está dentro da janela permitida ($8\text{h}$ às $18\text{h}$);
            \item \textbf{OU} o alarme de emergência está ativado (\texttt{True/False}).
        \end{itemize}
        Construa a expressão lógica em Python que determine se a porta deve destravar.
    \end{block}

    \begin{exampleblock}{Ajuda de Resolução (Como Pensar)}
        \begin{itemize}
            \item Identifique: \texttt{tem\_cracha} (bool), \texttt{hora} (int) e \texttt{emergencia} (bool).
            \item A janela de horário é: \texttt{(hora >= 8 and hora <= 18)}.
            \item Conecte as condições com os operadores \texttt{and} e \texttt{or}.
            \item Use parênteses para garantir que o \texttt{or} avalie a emergência como exceção geral.
        \end{itemize}
    \end{exampleblock}
\end{frame}

% Questao 5
\begin{frame}{Questão 05: Expressões Aritmético-Lógicas}
    \begin{block}{Quesito: Verificador de Regra de Ano Bissexto}
        Um ano é bissexto se atender às regras do calendário gregoriano:
        \begin{itemize}
            \item É divisível por $4$ \textbf{e} não é divisível por $100$; \textbf{OU}
            \item É divisível por $400$.
        \end{itemize}
        Solicite um ano inteiro e avalie se ele é bissexto usando \textbf{uma única expressão booleana} (sem \texttt{if/else}). Imprima apenas o resultado booleano.
    \end{block}

    \begin{exampleblock}{Ajuda de Resolução (Como Pensar)}
        \begin{itemize}
            \item Divisibilidade por $N$ é avaliada com operador de resto: \texttt{ano \% N == 0}.
            \item Não ser divisível por $100$: \texttt{ano \% 100 != 0}.
            \item Combine as condições em uma só linha usando operadores \texttt{and} e \texttt{or} com parênteses adequados.
        \end{itemize}
    \end{exampleblock}
\end{frame}
